\section{The platform}
\label{sec:experiments}

Section~\ref{sec:emptycell} ends at a cell nobody has filled. This section
describes the instrument we built to fill it. It is one harness. It runs a
published defense against a published attack under one protocol, with the
controls fixed before any run. The next two sections report what it
produced, and together they answer RQ2.

\subsection{Pair selection}

The taxonomy picks the (defense, attack) pair under two conditions. The defense and
the attack must occupy the same taxonomy cell, which means the content enters
through the same channel and causes the same consequence. And no paper in the
corpus may report having run them together. A defense that repairs quiet answer corruption arriving through retrieval
ought to apply to an attack that arrives the same way and does the same
damage. The taxonomy says so. Nobody had checked.

We are not proposing new defenses. Every method we run is published, and
every attack we run is published. The contribution is the pairing.

\subsection{The reconstruction standard}

The taxonomy decides which pairs we run, and this rule decides what counts as
running one. Every number we report comes from running a defense's own released code and
weights. We never cite a paper's reported numbers as evidence about that
paper's method. Where a harness had to be adapted to our data, we adapted the
harness and left the method and the evaluation rubric untouched. Where an
author's rubric and ours differed, we used the author's.

The standard cost us coverage. Several defenses could not be evaluated at
all, for reasons that lie in what their artifacts ship rather than in what
their methods do, and we record those as roadblocks rather than as failures
of the method.

\subsection{Pre-registered experimental rules}

Released code is not enough on its own. A negative result is easy to produce
by accident, so four rules guard against that, and we fixed all four before
any run.

\textbf{The control fires first.} A defense that is broken in our hands
tells us nothing about transfer, so before any transfer number we check that
the method still does, in a clean context, what its authors say it does. The check is directional, and it is set per method, because no two papers
claim the same quantity. ParamMute~\cite{huang-2025-parammute} reports a gain in context
recall, not in accuracy. CK-PLUG~\cite{bi-2025-parameters} claims no accuracy gain at all: its
contribution is a knob that moves memorization rate while accuracy is held,
so its control fires on that rate instead. ACD~\cite{kim-2024-adaptive} states its improvement
without a figure. SCR's~\cite{huang-2024-trust} own authors report it underperforming their
baseline on Llama-3-8B, the model we run it on, so on this model there is no
published gain to reproduce, and agreement with them means moving the same
way they report. For each method, we therefore check the direction of the change, on the metric
that paper uses for the model it used. We record which quantity each
control was read on. Where no comparable published figure
exists, we say so rather than manufacture a threshold.


\textbf{The attack works first.} If the attack does not degrade the undefended
arm in the same run, the defended arm means nothing. Every result below is
reported next to an undefended arm measured on the same items at the same doses
with the same scorer.

\textbf{Effects are gated on a minimum detectable change.} For each run we
compute the smallest difference detectable at the given sample size with
$\alpha = .05$ and power $.80$. Differences below it are reported as "no
resolvable effect", never as "no effect". This matters because several
published claims in this area rest on differences that a run of their size
cannot resolve.

\textbf{Attacks are not adapted to the defenses they meet.} An attack that is not built for a defense and
defeats it anyway is a lower bound on attack strength: if the
cheapest payload wins, an adapted one also wins. So non-adaptive evaluation is
conservative for a claim that a defense fails, and invalid for a claim that one
holds. Every failure we report is of the first kind. The single resistance
result we report is of the second, and we test it adaptively rather than
asserting it.

Predictions were written into the runner scripts and committed to version
control \emph{before} each run, and they ship with the artifact.

\subsection{What we ran}
\label{sec:whatweran}

Every pair we ran is two choices: which literature the defense came from,
and which literature the failure came from. That makes four boxes, and they
do not ask the same question.

Two boxes are home ground. A method meets the kind of failure its own authors
built it for, so nothing there tests transfer. These are where the first rule
above is satisfied, and they hold \NCampSecDefSecMech{} and
\NCampSafDefSafetyMech{} of our \NCampaignPairs{} pairs.

The other two boxes are crossings, and each asks its own question.

The first crossing puts an adversary in front of a method built for honest
error. These are methods for a context that is merely wrong --- out of date,
or contradicted by what the model already believes --- and the question is
what they do when someone chose the wrong text on purpose. That is most of
what we ran, \NCampSafDefSecMech{} pairs.

The second crossing runs the other way. It gives a method built against an
adversary a failure nobody wrote. The question is whether such a method needs
an attacker to earn its place, or whether it helps whenever the context stops
supporting the right answer. This is the smaller side,
\NCampSecDefSafetyMech{} pairs.

Inside each crossing the pairs vary along one axis, and it is a different
axis in each. Going one way, what changes is how much the attacker controls:
a passage optimized against the retriever, a fluent sentence asserting a
chosen answer, a passage that steers behavior rather than the answer, and a
passage optimized against the defense itself. Going the other way, what
changes is whether anything arrives at all: text that reaches the window and
should not be believed although no one wrote it to deceive, against a window
whose contents all deserve belief and which stops working because there is
simply more of everything.

The four boxes are not spread evenly across the two rigs the campaign runs
on. The retrieval cell holds all four. The agent cell holds one: every
technique in it is a prompt-injection defense and every attack an AgentDojo
injection, so all \NAgentPairs{} of its pairs are a security defense meeting
a security mechanism. Every crossing we report is therefore a retrieval-cell
result, and the agent cell tests no transfer at all.

The two crossings together are \NCampCrossPairs{} of \NCampaignPairs{}
pairs, \PctCampCross{}\% against the \PctCrossCommunity{}\% the
literature manages,
and \NRagImplSafety{} of our \NRagImpl{} retrieval defense implementations
are knowledge-conflict methods meeting an attack. That is the literature's
empty column, populated. Counted at registry-name level,
\NRagPairsNovel{} of \NRagPairsKeyed{} pairs are absent from the
\NPairsComposite{} the literature reports, and the two that are not are the
home-ground pairs their own papers already ran.

The retrieval cell pairs defenses built for knowledge conflict against
attacks built to poison a retrieval corpus. The victim model is
Llama-3-8B-Instruct~\cite{grattafiori-2024-llama3} or Mistral-7B-Instruct~\cite{jiang-2023-mistral}, depending on what each defense
released. The data is CoConflictQA~\cite{huang-2025-parammute} over Natural Questions, and open Natural
Questions for the attacks that need a retrieval pool.

Those defenses turn the same dial. ACD decodes from the no-context
distribution and pulls it toward the with-context one by a weight computed
from entropy. CK-PLUG interpolates the two
distributions under a single parameter. ParamMute
suppresses the feed-forward layers that carry parametric memory, forcing the
model onto the retrieved text. SpARE~\cite{zhao-2024-steering} and
SHIFT~\cite{li-2026-shift} steer along a context-faithfulness direction in
activation space. COIECD detects a conflict
between the two sources and pushes toward the context.
SCR makes the same trade through prompting rather
than decoding. Each changes \emph{how much} the model leans on retrieved
text. None asks \emph{which} retrieved text deserves trust.

The agent cell pairs ingestion-stage and execution-stage defense
techniques against prompt-injection attacks inside the AgentDojo banking
suite~\cite{debenedetti-2024-agentdojo}, with a Qwen2.5-14B-Instruct~\cite{qwen-2024-qwen25} agent.

In total we executed \NCampaignPairs{} distinct (defense, attack) pairs
across \NCampaignImpl{} defense implementations, over \NCampaignEpisodes{}
scored episodes: \NRagPairs{} pairs and \NRagEpisodes{}
episodes in the retrieval cell, \NAgentPairs{} pairs and \NAgentEpisodes{}
episodes in the agent cell.
The two episode counts are not comparable units. An agent episode is a complete
multi-step tool-use trajectory. A retrieval episode is one scored answer.

Every run in this paper used one NVIDIA RTX A6000 with 48\,GB, driver
535.274.02 and CUDA 12.8, under Python 3.10.12, PyTorch 2.10.0 and
\texttt{transformers} 5.5.0. Two defenses cannot use that last version. ACD's
released decoder imports symbols that \texttt{transformers} removed after
4.44, which its own README pins, so it runs in a separate dependency set.
CK-PLUG ships its own 4.49 fork and its README says to install it. We record this
because it is a result as much as a configuration detail. The defenses in a
single cell of this taxonomy cannot share one Python environment. That is the
reproducibility problem of Section~\ref{sec:discussion}, one layer further
down. The released artifact builds each set and verifies it.

\subsection{Datasets}
\label{sec:provenance}

An \textbf{item} is one question, its gold answer, and the context the model
is given to answer it from. Every item in this paper belongs to the authors
of the method or the benchmark we are running. Table~\ref{tab:provenance} lists each source and whose it is. This is deliberate rather than convenient. A defense
evaluated on the distribution it was designed for cannot blame its failure on
an unfamiliar dataset. A result obtained that way is harder to dismiss than
one obtained on data we chose ourselves.

\begin{table}[t]
\centering
\caption{Item sets and payloads, by owner. No item set is ours. Three
payloads come from the attack literature and five are ours, and none of the
five is a new attack: three are edits of a field the benchmark already ships,
and two reuse the item set's own retrieved passages, which is what makes them
controls rather than attacks.}
\label{tab:provenance}
\small
\begin{tabular}{@{}>{\raggedright\arraybackslash}p{0.40\columnwidth}>{\raggedright\arraybackslash}p{0.52\columnwidth}@{}}
\toprule
\textbf{Item set} & \textbf{Released by} \\
\midrule
CoConflictQA~\cite{huang-2025-parammute} over Natural Questions & the ParamMute authors \\
Open Natural Questions~\cite{kwiatkowski-2019-natural}, with its retrieval pool & the RobustRAG authors \\
FaithEval~\cite{ming-2025-faitheval}, SQuAD~\cite{rajpurkar-2016-squad} negatives & the FaithfulRAG~\cite{zhang-2025-faithfulrag} authors \\
AgentDojo banking suite & the AgentDojo authors \\
\midrule
\textbf{Payload} & \textbf{Origin} \\
\midrule
Corpus poisoning & Zhong et al.'s released passage, unmodified \\
PoisonedRAG, BadRAG & The published constructions \\
False evidence & CoConflictQA's own parametric-evidence field \\
Distractor evidence & \textbf{ours}: one substitution into that field \\
Adaptive distractor & \textbf{ours}: the same claim, reshaped against one
defense's rule \\
Gradient distractor & \textbf{ours}: that passage and a suffix optimized
against the same rule \\
Stale document & \textbf{ours}: a real retrieved passage, gold answer
swapped \\
Misretrieval & \textbf{ours}: a real passage, unmodified, retrieved for
another question \\
\bottomrule
\end{tabular}
\end{table}

Every comparison in this paper is between arms measured on the same items.
Rates computed on different item sets are incomparable rather than merely
noisy, so we checked this mechanically rather than by inspection: across
every result registry, no run compares two arms measured on different items.

The one place the item set narrows is the distractor payload, and it narrows
for both arms together. It is defined on \NDistractorItems{} of CoConflictQA's
\NCoConflictItems{} items. An item
qualifies on two conditions. Its remembered answer must appear verbatim in its
own evidence sentence, so that the substitution needs no rewriting. And its
question type and answer magnitude must admit at least three candidate
distractors, so that the planted answer is a plausible example of its kind
rather than an absurdity a reader would dismiss. Every arm in those runs, including the undefended and no-retrieval
arms, is measured on that same subset.

\subsection{Payload classes}
\label{sec:payloadclasses}

The retrieval cell evaluates every item whose context is at most 2000
characters. The cut is on context length alone, never on an outcome, and it
is applied once and then frozen, so every defense and every payload meets the
same items. Its purpose is memory. Contrastive-decoding defenses hold two
forward passes of the model at once, and this keeps the whole cell inside one
card at a batch size every defense can share.

On those items we hold the dosing and the scorer fixed as well, and vary only
the \emph{form} of the planted passage. There are five forms. They differ in
how the passage reads and in what it does to the evidence, and a defense that
behaves differently across them tells us which property it is actually keying
on.

Corpus poisoning~\cite{zhong-2023-poisoning,su-2024-corpus} optimizes a
token sequence to be retrieved for many queries. The resulting passage is
incoherent to a human reader. It attacks by displacing honest documents.

False evidence does the opposite. It plants a fluent, well-formed sentence
that asserts the wrong answer the model itself would have given from memory. It is drawn from
CoConflictQA and reads like an ordinary encyclopaedia sentence. It attacks by
agreeing with the model's prior.

Distractor evidence is that same fluent sentence with one word changed, so
that it asserts a third answer which is neither the truth nor the model's
prior. It is not a stronger attack. It is a measuring instrument, and
Section~\ref{sec:distractor} explains why it was necessary.

An adaptive distractor is the distractor payload again, in whichever surface
form minimizes the quantity one defense's rule measures. A gradient
distractor is that passage with a short optimized suffix appended against the
same quantity (Section~\ref{sec:results2}).

Three of the five are ours, the distractor and its two optimized forms, and
all three are edits of a field that already ships with the benchmark. That field is worth
drawing out. False evidence is the
\texttt{parametric\_evidence} field that ships inside CoConflictQA, the
benchmark the ParamMute authors released themselves. One of the defenses we
break is broken with a field from its own dataset.

\paragraph{Dose, and the one item set that has none.}
The \textbf{dose} is how many of the retrieved passages are planted. It is
the axis every result in this paper is read along, because a defense that
holds against one poisoned passage and fails against five has not failed in
the same way as one that never helped.

We take the dose from the attacks rather than define one. Their released
constructor is parameterised by how many of the top $k$ retrieved passages
it replaces, where it places them and how often the planted sentence repeats
inside its passage. We keep that construction and sweep the one parameter:
$j$ of the top ten, for $j$ in $\{0, 1, 2, 3, 5, 10\}$. Their own default is
$j = 1$, so their setting is one of our points and the rest are further
readings of the same instrument.

One item set has no such axis and we do not give it one. CoConflictQA
supplies a single context passage per item, which either carries the
conflicting claim or does not, and a single passage has no ranked list to
dose. We report it in the two conditions its own authors run, a clean
context and a corrupted one. The two item sets are never combined in one
number.

Every result file also carries a fingerprint of the payload it was produced
with, so a run cannot be read under the wrong attack label.
